% Einheit 3 von 3 – Kugel (bank/pyramide-kegel-kugel/e3.jsonl, zusammenbau v0.1)
%% TODO Zweigzeile Teil 1: was hier gelernt wird (Satz in Schülersprache aus dem Einheitstitel) – Einheit 3
\einheitenkopf[e3]{Einheit 3 von 3 · Kugel}
\zweigzeile{ab Kl. 8, je nach Buch bis Kl. 10 (am Gymnasium ab Kl. 9) · P10}
\setcounter{aufgabe}{20}

%% TODO Titel: Ich-kann-Satz fehlt in der Bank (Platzhalter: Kettenname) – Nr. 21
%% TODO Anweisung: ein Satz über der ersten Teilaufgabe fehlt in der Bank – Nr. 21
\begin{aufgabe}{Kugel}
%% TODO \rechenplatz{n}: Schrittzahl des Grundfalls fehlt in der Bank (nur bei fester Schreibform, 2.3 b)
\begin{teile}
\teil Die Strecke mit ? an der Kugel – Radius oder Durchmesser? Kreuze an.\\ \kreuz{Radius}\\ \kreuz{Durchmesser}

\kugel{1.5}{?}
\teil Kugel: Radius $3$ cm – Volumen? Runde auf eine Stelle. $V$ ≈ \leerfeld[cm³]
\teil Kugel: Durchmesser $16$ cm – Volumen? Runde auf eine Stelle. $V$ ≈ \leerfeld[cm³]
\teil Kugel: Radius $2{,}5$ cm – Volumen? Runde auf eine Stelle. $V$ ≈ \leerfeld[cm³]
\teil Ein Wasserball hat den Radius $15$ cm. Wie viel Liter Luft sind darin? \leerfeld[l]
\teil Kugel: Radius $11$ cm – Oberfläche? Runde auf eine Stelle. $O$ ≈ \leerfeld[cm²]
\teil Skizziere eine Kugel mit dem Radius $2$ cm: Kreis, Äquator als Ellipse, Radius gestrichelt und beschriftet.

\rechenplatz[halb]{4}
\teil Halbkugel: Radius $4$ cm – Volumen und Oberfläche? Runde auf eine Stelle. $V$ ≈ \leerfeld[cm³] ; $O$ ≈ \leerfeld[cm²]
\teil Eine Stahlkugel hat den Radius $2$ cm, Stahl hat die Dichte $7{,}9$ g/cm³. Wie schwer ist sie? Runde auf ganze Gramm. $m$ ≈ \leerfeld[g]
\teil Kugel: Oberfläche $200$ cm² – Radius? Runde auf zwei Stellen. $r$ ≈ \leerfeld[cm]
\teil Ein Ball mit dem Radius $14$ cm soll genau in eine würfelförmige Schachtel passen. Wie lang ist die Kante? Wie viel cm³ Luft bleiben in der Schachtel? \leerfeld[cm] ; \leerfeld[cm³]
\teil Ein Silo besteht aus einem Zylinder (Radius $3$ m, Höhe $9$ m) mit einem Halbkugeldach. Wie viel m³ fasst es insgesamt? \leerfeld[m³]
\teil Kugel A hat den Radius $2$ cm, Kugel B den Radius $4$ cm. Berechne beide Volumen und zeige, dass B achtmal so viel Volumen hat.
\teil Eine Bowlingkugel hat den Durchmesser $21{,}6$ cm. Berechne ihr Volumen mit der Formel aus der Formelsammlung \hfill (P10 2016 OS). $V$ ≈ \leerfeld[cm³]
\end{teile}
\end{aufgabe}

%% TODO Titel: Ich-kann-Satz fehlt in der Bank (Platzhalter: Kettenname) – Nr. 22
%% TODO Anweisung: ein Satz über der ersten Teilaufgabe fehlt in der Bank – Nr. 22
\begin{aufgabe}{Kugel – Fehler finden, Begründen, Darstellungswechsel, Anwendung}
\begin{teile}
\teil Eine Kugel hat den Durchmesser $18$ cm. Nora rechnet das Volumen: \rechnung{V &= \tfrac{4}{3} \cdot \pi \cdot 18^3 \\ V &\approx 24\,429{,}0 \text{ cm}^3} Wo ist der Fehler?
\teil Warum ist die Oberfläche einer Halbkugel mehr als halb so groß wie die Oberfläche der ganzen Kugel? Begründe.
\teil Ein Silo besteht aus einem Zylinder mit dem Radius $2{,}5$ m und der Höhe $8$ m, oben sitzt ein Halbkugeldach. Skizziere das Silo und beschrifte Radius, Höhe des Zylinders und Gesamthöhe.

\rechenplatz[halb]{4}
\teil Die Kuppel eines Turms ist eine Halbkugel mit dem Radius $2{,}8$ m. Sie wird außen zweimal gestrichen; ein Eimer Farbe reicht für $15$ m². Wie viele Eimer braucht man?
\end{teile}
\end{aufgabe}
